\section{Modelagem matemática}
\label{sec:modelagem-matematica}

Ambos os domínios, $\Omega_f$ e $\Omega_p$, são definidas pelo mesmo conjunto de equações, o que os distingue é a presença ou não da interface e, por consequência, o tratamento dado às propriedades do fluido. A formulação é apresentada aqui em sua forma mais geral, bifásica, da qual o escoamento resolvido em $\Omega_p$ é recuperado como caso particular monofásico.

Adota-se a formulação de um fluido (\textit{one-fluid formulation}), na qual as duas fases são descritas por um único conjunto de equações válido em todo o domínio. Em lugar de resolver um sistema de equações por fase, acoplados por condições de salto na interface, a massa específica $\rho(\vec{x},t)$ e a viscosidade dinâmica $\mu(\vec{x},t)$ passam a ser campos descontínuos através da interface.

Sendo os campos de propriedades variáveis, o balaço de massa é tomado em sua forma geral, Eq. \ref{eq:conservacao-massa}:

\begin{equation}
    \frac{\partial \rho}{\partial t} + \nabla \cdot \left( \rho \vec{u} \right) = 0
    \qquad \Longleftrightarrow \qquad
    \frac{D \rho}{D t} + \rho \, \nabla \cdot \vec{u} = 0
    \label{eq:conservacao-massa}
\end{equation}

Sendo os fluidos são incompressíveis e imiscíveis e não havendo mudança de fase. Cada partícula conserva, portanto, sua massa específica ao longo do escoamento, de modo que $D\rho / Dt = 0$, e a interface é uma superfície material, transportada pelo próprio campo de velocidade e sem fonte de volume associada. O primeiro termo da Eq. \ref{eq:conservacao-massa} se anula e resta a condição de divergência nula do campo de velocidade, Eq. \ref{eq:continuidade-completa}, válida mesmo com $\rho$ variável no domínio: a variação da massa específica decorre da distribuição das fases, e não da compressão do fluido. 

\begin{equation}
    \nabla \cdot \vec{u} = 0
    \label{eq:continuidade-completa}
\end{equation}

O balanço de quantidade de movimento é dado pela Eq. \ref{eq:navier-stokes-incompressivel-completa}, com a dependência espacial e temporal das propriedades explicitada:

\begin{equation}
    \frac{\partial \vec{u}}{\partial t} + \vec{u} \cdot \nabla \vec{u}
    = \frac{1}{\rho (\vec{x},t)} \left[
        - \nabla p
        + \nabla \cdot \left( \mu (\vec{x},t) \left( \nabla \vec{u} + (\nabla \vec{u})^{T} \right) \right)
        + \vec{f}_{\sigma}
        + \vec{f}_{ib}
    \right]
    + \vec{g}
    \label{eq:navier-stokes-incompressivel-completa}
\end{equation}

nas quais $\vec{u}$ é o vetor velocidade, $p$ a pressão, $\vec{g}$ o vetor aceleração da gravidade, $\vec{f}_{\sigma}$ a força por unidade de volume devida à tensão superficial, não nula apenas sobre a interface, e $\vec{f}_{ib}$ a força por unidade de volume responsável por impor a condição de contorno sobre a fronteira imersa, quando presente.

As particularidades na implementação de cada fluido pertencem ao modelo numérico e serão tratadas de forma separada.

